\chapter{沉淀、配位与化学侦探术}

\begin{kaichang}
找一个用了半年以上的电热水壶，打开盖子往里看：壶底和壶壁上多半结着一层灰白色的硬壳，用手指抠一抠，糙得像砂纸。这就是水垢。很多人家的处理办法是倒点醋进去泡一泡——你会看到白色硬壳表面冒出细密的小气泡，慢慢地，硬壳变薄、消失。

先别急着觉得这件事平淡无奇。请认真想一想：自来水看起来清澈透明，烧开了也干干净净，\textbf{这层白色的东西是从哪儿冒出来的？}如果水里真的藏着“料”，为什么我们看不见？还有，醋明明没有“磨”它，只是泡着，硬壳怎么就消失了，还冒泡？

再追问一步：自来水厂、矿泉水公司、游泳池管理员，都声称自己“检测”过水质。水里的东西看不见摸不着，他们是怎么知道里面有什么、有多少的？先猜一个答案。这一章结束时，你不仅能解释水垢，还会拥有化学家当“侦探”的三件法宝。
\end{kaichang}

\section{水壶里的“隐形客人”}

先把能看见的事实摆出来。

事实一：清澈的水烧久了，会结出白色固体。这说明水里原本溶解着某些\textbf{看不见的物质}，烧水这个过程把它们“逼”了出来。事实二：这层白色固体怕醋，泡在醋里会冒气泡、逐渐溶解。事实三：如果你家有桶装纯净水，用同样的壶只烧纯净水，水垢几乎不出现。可见水垢的原料不是“水本身”，而是水里的“乘客”。

类似的现象生活里到处都是。烧水壶的壶嘴内侧、淋浴喷头的小孔、暖水瓶的内胆，时间一长都会挂上白霜一样的东西。钟乳石洞里更壮观：从洞顶垂下的石笋、石柱，也是水一滴一滴“送”来的白色固体，只不过大自然用了几十万年。还有一种更隐蔽的：用硬水（含钙镁较多的水）洗衣服，肥皂好像变“懒”了，泡沫少，水面上还漂着一层灰白色的浮渣——那也是水里的隐形客人在作怪。

这些客人是谁？要回答这个问题，得把镜头推进到粒子层。

\section{客人现形：离子在溶液里游荡}

第 25 章已经讲过，溶液里发生的事远比肉眼看到的精彩：食盐溶于水后，并不是“盐消失了”，而是氯离子和钠离子被水分子团团围住、各自散开游荡（第 23 章讲过水分子怎样用它的极性“拆散”离子晶体）。河水、井水、自来水在地下流过岩石和土壤时，同样把石灰岩（主要成分碳酸钙）、石膏等矿物里的离子“请”进了水里——其中最多的客人就是钙离子 \chem{Ca^{2+}} 和碳酸氢根离子 \chem{HCO_3^-}。

这些离子平时相安无事，因为它们在水里跑得太欢，浓度又不高。可一旦烧水，情况变了：温度升高，碳酸氢根变得不稳定，会转变成碳酸根 \chem{CO_3^{2-}} 并放出二氧化碳。碳酸根和钙离子一旦相遇，彼此间的静电吸引力很强（离子键的吸引力，第 18 章讲过），结合成的碳酸钙又很难被水分子“拆开带走”，于是它们一批批地从水里析出、堆积成固体——这就是水垢，主要成分就是石灰岩的那位亲戚：碳酸钙 \chem{CaCO_3}。

所以钟乳石和水垢本质上是同一出戏，只是舞台不同：溶有钙离子和碳酸氢根的水从洞顶渗出来，水分蒸发、二氧化碳逸出，碳酸钙就一点点析出，几十万年垒出一根石柱。

那醋为什么能除掉水垢？醋里的醋酸提供氢离子 \chem{H^+}（第 32 章讲过酸的本质就是给质子），氢离子抓住碳酸根，把它转变成二氧化碳气体和水——你看到的细密气泡就是二氧化碳。固体被“拆”成能溶于水的离子和跑掉的气体，硬壳自然就消失了。沉淀能被造出来，也能被拆掉，这正是本章第一件法宝的雏形。

\section{多浓才算“太挤”：溶度积}

“难溶”这个词太含糊了。碳酸钙难溶，到底是多难？化学家需要一把定量的尺子。

先把一个重要观念搬过来：第 31 章讲过，化学平衡不是停住不动。一块碳酸钙固体泡在水里，表面上风平浪静，粒子层面却两班倒：固体表面的钙离子和碳酸根不断脱落进水，水里的钙离子和碳酸根也不断撞回固体表面重新结晶。两者速率相等时，溶解和析出达到动态平衡，溶液就“饱和”了。

\begin{qiaoliang}
对碳酸钙这样的难溶盐，平衡时溶液里两种离子浓度的乘积是一个固定值（温度不变时），叫做\textbf{溶度积}，记作 $K_{\mathrm{sp}}$：
\[K_{\mathrm{sp}}(\chem{CaCO_3}) = [\chem{Ca^{2+}}]\times[\chem{CO_3^{2-}}] \approx 3\times10^{-9}\quad(25\,^{\circ}\mathrm{C})\]
判断方法像查账：把此刻两种离子的实际浓度相乘，记作 $Q$。$Q$ 小于 $K_{\mathrm{sp}}$，说明溶液还“装得下”，固体会继续溶解；$Q$ 一旦超过 $K_{\mathrm{sp}}$，溶液就“挤爆了”，多余的离子只好析出成沉淀。沉淀不是“反应到底”，而是被这把尺子拦在了门外。
\end{qiaoliang}

拿这把尺子算一笔账。设碳酸钙在纯水中的溶解度为 $s$，平衡时两种离子浓度都等于 $s$，所以 $s^2\approx3\times10^{-9}$，开平方得 $s\approx5\times10^{-5}\ \mathrm{mol/L}$。碳酸钙的摩尔质量约 \ud{100}{g/mol}，折算下来每升水大约只能溶解 \ud{5}{mg}——一杯水里的溶解量还不到一粒盐的重量。“难溶”二字，就这样变成了具体数字。

再算一笔和生活有关的账。很多城市的自来水硬度以碳酸钙计约 \ud{200}{mg/L}。假设你家的壶 \ud{1.5}{L}，每天烧一壶，一年流过壶的水约 $1.5\times365\approx550$ 升，携带的“水垢原料”约 $200\times550\approx1.1\times10^{5}\ \mathrm{mg}$，也就是一百多克。哪怕只有三分之一真正结成垢，一年也有三十多克——够把壶底糊上一层硬壳了。你看，一个让人皱眉的生活麻烦，用离子浓度和一个乘积就算得明明白白。

沉淀反应在医学上也有一位大明星。做胃肠造影检查时，医生会让病人喝下一杯白色的“钡餐”——硫酸钡 \chem{BaSO_4} 悬浊液。这里有个看似矛盾的地方：可溶性的钡盐是剧毒，钡离子会干扰神经和肌肉，可钡餐天天在喝。秘密就在溶度积：硫酸钡的 $K_{\mathrm{sp}}$ 只有约 $1\times10^{-10}$，溶解度低到每升水只溶出约两毫克，喝下去几乎不被肠道吸收，原样穿肠而过，又能挡住 X 射线显出肠胃轮廓。毒性从不只看“是什么元素”，更要看它以什么形态存在、能溶出多少——第 33 章讲氧化态时我们已经见过这个道理。

\section{侦探的第一招：让离子自己“现形”}

现在回到开头的大问题：怎样知道水里“有什么”？化学侦探的第一招，是利用沉淀反应的\textbf{专一性}——让目标离子变成一件看得见的“证物”。

最经典的案例是检验氯离子。向水样里滴加硝酸银溶液，如果出现白色沉淀，嫌疑犯就是氯离子，因为生成的氯化银 \chem{AgCl} 极难溶（$K_{\mathrm{sp}}\approx1.8\times10^{-10}$），立刻析出。用符号写，就是把看不见的过程记成一行账：
\[\chem{Ag^+} + \chem{Cl^-} \rightarrow \chem{AgCl}\downarrow\]
同样地，检验硫酸根 \chem{SO_4^{2-}} 用氯化钡溶液，白色的硫酸钡沉淀就是“证物”；检验水样硬度，可以利用钙镁离子让肥皂水起浮渣的特性。

但真正的侦探都知道：\textbf{出现一个证据，不等于锁定了凶手。}硝酸银遇到碳酸根也会生成沉淀，如果你不问青红皂白就下结论，就可能“冤判”。所以标准的氯离子检验流程里有一步看似多余的操作：先加稀硝酸酸化。碳酸银等“冒牌沉淀”会溶于酸（碳酸根被氢离子变成二氧化碳跑掉），而氯化银不溶于酸，留到最后。排除干扰，是分析化学里和“产生信号”同样重要的一半功夫。侦探的思维不是“有反应就下结论”，而是“还有什么别的东西也能给出同样的现象？我怎样把它排除掉？”

\begin{shiyan}
\textbf{用肥皂水给水“验明正身”}

材料：两个一样的透明杯子、一瓶纯净水、自来水（或矿泉水）、液体肥皂或洗洁精、一根筷子。

步骤：两个杯子各装半杯水，一杯纯净水、一杯自来水。各滴入同样多的肥皂液（比如各 10 滴），用筷子以同样的速度搅同样长的时间，或者盖紧杯盖摇 10 下，然后静置观察。

观察点：哪一杯泡沫丰富持久？哪一杯泡沫稀少，水面漂着灰白色的浮渣？浮渣就是钙、镁离子和肥皂分子结合生成的不溶物——泡沫越“懒”，说明水里的钙镁离子越多，水越“硬”。

安全提示：实验用的水\textbf{不要再喝}；肥皂液别弄进眼睛；结束后把杯子洗净。这个实验只用家常物品，可以放心做。
\end{shiyan}

\section{侦探的第二招：给离子“穿外套”}

第一招是把离子“抓出来”，第二招恰好相反：给它“穿上外套”，让它变得更显眼或更安分。

在第 32 章，我们把酸碱的定义从“给出质子”扩展到了“接受或给出电子对”。沿着这个思路走一步：金属阳离子是电子对的“空房子”，而有些分子或离子带着富余的电子对，像热情的租客。当它们靠上去、用电子对和中心离子配成键，就形成一个整体，叫做\textbf{配离子}（也叫络离子）；提供电子对的分子或离子叫\textbf{配体}，被围在中间的正离子叫\textbf{中心离子}，直接贴在中心离子周围的配位原子个数叫\textbf{配位数}。这种“租客—房东”式的结合，就是配位键——本质上是共价键的一种，只是这对电子全由一方提供（第 19 章讲过共价键）。

最有名的演示：浅蓝色的硫酸铜溶液里，铜离子其实是被六个水分子围着的配离子，配位数是 6。滴加少量氨水，先出现浅蓝色沉淀（氢氧化铜）；继续加过量氨水，沉淀竟然溶解了，溶液变成深得发亮的宝蓝色——氨分子挤走了水分子，生成四氨合铜离子，中心离子 \chem{Cu^{2+}} 被四个氨分子抱住，配位数是 4：
\[\chem{Cu^{2+}} + 4\chem{NH_3} \rightarrow [\chem{Cu(NH_3)_4}]^{2+}\]
这个颜色变化本身就是一个侦探信号：那抹标志性的深蓝，是“溶液里有铜离子”的招牌证词。第 15 章讲过渡金属时我们见过它们的五颜六色，那些颜色大多来自配离子——中心离子穿了什么“外套”，颜色就不一样。

配位不只是颜色游戏，它有实打实的用途。洗衣粉里常加一类能“抱住”钙镁离子的配位剂，让硬水中的钙镁不再捣乱；锅炉和水管工人用类似的物质溶解水垢——把水垢里的钙离子“拽”回溶液。医学上，重金属中毒（如铅中毒）的一种急救手段是“螯合治疗”：注射像 EDTA 这样的药物，它像螃蟹钳子一样用六个配位原子把铅离子牢牢夹住，形成稳定又易溶于水的配离子，随尿液排出体外。“螯”就是螃蟹钳的意思——六个钳脚同时夹住一个离子，比一个钳脚牢固得多。当然，这类药物本身也会夹走身体需要的钙、锌，必须由医生严格控制剂量和疗程，绝不能自行使用。

现在，水垢那笔账也能重新理解了：碳酸钙的“难溶”不是绝对的。酸能用氢离子拆掉碳酸根，配位剂能用电子对拽走钙离子——沉淀、溶解、配位，三者是同一盘平衡棋局里的不同走法。

\section{侦探的第三招：看颜色、称质量}

前两招够聪明，但它们有个局限：得预先猜“凶手可能是谁”，才能选对试剂。如果面对一个完全未知的样品，或者目标物质含量低到十万分之一，试剂法就力不从心了。现代化学侦探靠的是三件仪器法宝，它们分工明确。

第一件法宝是\textbf{光谱}：让物质自己“报名字”。把样品烧进火焰或电弧里，原子被激发后会放出特定颜色的光，这光经过棱镜分开，是一组特定位置的亮线——像条形码一样，每种元素的“条码”全世界独一无二，不会和任何别的元素重复。在样品里看到钠的黄色亮线，就知道钠在场；看亮线的强弱，还能估计含量。原子吸收光谱反过来用：让特定元素的光穿过样品蒸气，被吸收多少，就说明有多少这种原子。水里的铅、镉、汞，土壤里的重金属，都是这样测出来的，灵敏度可以到十亿分之一量级。

第二件法宝是\textbf{色谱}：负责把混合物\textbf{分开}。你可能在生物课上做过：把菠菜叶的色素提取液点在滤纸条上，纸条下端浸进溶剂，溶剂往上爬，各种色素因为“爱溶剂”和“爱纸”的程度不同，爬得有快有慢，最后分成几条色带——绿色的叶绿素、黄色的叶黄素、橙色的胡萝卜素各就各位。现代色谱把这个思路做到极致：一根细管子里，几十种成分在几分钟内鱼贯而出，每个成分出场的“时间点”就是它的身份证。运动员的尿检、食品里的农药残留筛查，第一步几乎都是色谱。

第三件法宝是\textbf{质谱}：负责给分子\textbf{称体重}。把分子打成带电碎片，让它们在电场磁场里飞行，越重的粒子转弯越“迟钝”，于是按质量排好队。称出精确的分子量和碎片质量，再对照数据库，分子的身份就八九不离十。机场安检口“擦一擦行李试纸”查爆炸物、警方鉴定毒品、科学家分析火星土壤，用的都是它。

实际办案时三件法宝常联手：色谱先把混合物分开，质谱和光谱再逐个认人——这就是新闻里“检出了某某物质”背后的流水线。

\begin{zhengju}
1860 年，德国化学家本生和物理学家基尔霍夫做了一件当时看来很奇怪的事：他们把几吨矿泉水蒸干，取残渣放进火焰烧，再用分光镜看火焰的光。为什么要费这么大劲？因为他们有一个清晰的问题：\textbf{元素的“光谱指纹”能不能用来发现未知元素？}

设计很讲究。本生特意改进了一种温度高、本身几乎无色的燃气灯（就是今天实验室里的本生灯），这样火焰不会用自己的颜色淹没样品的信号；把几吨水浓缩成一点残渣，是为了让微量成分的谱线强到能被看见。当他们观察残渣的光谱时，熟悉的线条之外出现了\textbf{两条从未见过的亮蓝色线}。

接下来是关键一步：排除替代解释。会不会是已知元素在特殊条件下产生的新线？他们把当时所有已知元素的光谱逐一比对，没有一种能对上。会不会是仪器假象？换装置、换光源、反复重测，两条蓝线始终在原位。结论只剩一个：这里有一种没人见过的元素。他们给它起名“铯”（\chem{Cs}），来自拉丁语“天蓝色”。第二年，他们又在矿石光谱里看到两条陌生的暗红线，发现了铷（\chem{Rb}，“深红色”）。更戏剧性的是 1868 年，天文学家在太阳光谱里发现一条对不上的黄线——那个元素干脆先在太阳上被“看见”、27 年后才在地球上找到，它叫氦。

这个故事里最值得学的不是“光谱真厉害”，而是方法链条：先有明确的问题（指纹能不能指认新元素），再有针对性的设计（无色火焰、浓缩样品），然后穷尽替代解释（比对全部已知谱线），最后才敢宣布“这是新东西”。侦探术的核心从来不是仪器，而是这套“让错误结论无处藏身”的思维。
\end{zhengju}

\section{侦探也会失手：这些方法的边界}

每一件法宝都有它的“使用说明书之外”。

溶度积模型在稀溶液里很好用，但浓度高了就会走样：真实溶液里离子互相牵扯、还有配体来“抢人”，有效的浓度（活度）并不等于你称出来的浓度。所以氯化银在纯水里难溶，在浓氨水里却能溶掉——氨分子把银离子配位“拽”走了，溶解平衡被整个搬动。“不溶”永远是有条件的不溶，温度变了、配体来了，结论就得重算。

试剂检验依赖“猜对嫌疑人”。它回答的是“我猜的那个东西在不在”，而不是“这里面到底有什么”。面对完全未知的样品，试剂盒会集体沉默，这时必须请出光谱、色谱、质谱。仪器也有底线：每种方法都有\textbf{检出限}——低于这个含量，信号就淹没在噪声里，只能说“没检出”，不能说“绝对没有”。“未检出”和“不存在”之间隔着一条灵敏度的鸿沟，读检测报告时千万别混淆。

还有样品本身的陷阱。检验一包奶粉之前，采样的那一勺能不能代表整批货？容器会不会把外面的铅带进样品？侦探小说里最精彩的是推理，真实实验室里最容易翻车的却是取样和污染。

\begin{wuqu}
“新闻说某食品里检出了某种化学成分，所以这食品有毒，不能吃。”——这是最常见的恐慌逻辑，它混淆了“检出”和“有害”。

仪器越灵敏，“检出”就越容易，而“有害”完全是另一回事。砒霜里的砷够可怕了吧？可每个人的血液、每一杯海水里都能检出微量砷；香蕉天生带一点放射性钾-40，照吃无害；连纯净水里都能检出痕量铅——关键是含量距离“可能造成危害的剂量”有多远。毒理学的老话说得直白：\textbf{剂量决定毒性}。连水喝得太多太快都会中毒，而剧毒的肉毒杆菌毒素稀释到亿万分之一，反而是医生手里治疗肌肉痉挛的药。

所以看到“检出”二字，该问的不是“怎么会有这种东西”，而是：检出了多少？安全限值是多少？超没超标？证据是单次检测还是反复验证？风险评估从来不是“有或没有”的是非题，而是“多少、多久、对谁”的算术题。这也是第 12 章以来我们反复强调的态度：化学名词不可怕，不讲剂量才可怕。
\end{wuqu}

\section{回到那只水壶}

现在可以完整地解释开场的一幕了。

自来水从岩层里一路走来，溶解了钙离子和碳酸氢根——它们就是看不见的“隐形客人”。水烧开时，碳酸氢根受热分解出碳酸根，水里钙离子和碳酸根的浓度乘积超过了碳酸钙的溶度积，析出的固体一层层糊在壶底，成了水垢。倒进醋以后，醋酸提供的氢离子把固体里的碳酸根“变”成二氧化碳气体（那些细密的气泡）和水，钙离子则留在溶液里——白色硬壳就这样被“拆”掉了。除垢剂厂商用的正是同一套思路，只不过把醋酸换成了更强的酸或更专业的配位剂。

至于自来水厂的化验员，他们的工作你也能想象了：用肥皂水或络合滴定测硬度（用 EDTA 溶液一滴滴“夹走”钙镁离子，指示剂变色那一刻读出总量），用硝酸银法或专门的电极测氯离子，用原子吸收光谱盯重金属，用色谱和质谱筛查农药和有机污染物。每一杯“合格”的自来水背后，都是一整套侦探流水线。

\section{第一册的终点，下一程的起点}

这本书的第一篇就叫“物质侦探的工具箱”。走到第一册的最后一章，你的工具箱已经装得相当满了：你知道物质由原子组成，原子按电子楼层排成周期表；原子靠离子键、共价键、金属键结伴，分子靠分子间作用力聚散；水用它的极性溶解万物，溶液里的反应受热、熵、速率、平衡、酸碱、氧化还原的共同指挥。这一章的沉淀、配位和仪器分析，是这套工具的一次“综合办案”。

但侦探故事还有一个更大的悬案没破。这一章我们一直在追问“水里有什么”，可你有没有想过追问“\textbf{你}是由什么搭起来的”？你的皮肤、头发、肌肉，你吃的米饭、鸡蛋、牛肉，你呼吸时交换的气体——这些分子大多不是本章这种小小的离子和沉淀，而是成百上千个原子搭成的复杂结构。法医能从一根头发里读出主人的身份，质谱能称出这些分子的体重，可要理解它们为什么长那样、为什么能干那些事，光靠无机离子的规则就不够了。

答案要交给一种特殊的原子。它能和自己连成链、绕成环、搭出枝杈，像乐高积木一样拼出几百万种不同的分子——它就是碳。第二册《碳怎样活了起来》将从第 36 章开始，看看碳凭什么成为分子世界的“搭骨架大师”，看这些碳骨架怎样一步步搭出糖、脂肪、蛋白质，最终搭出“活”这件事本身。无机世界的侦探术已经到手，接下来，我们去侦破生命。

\begin{tiaozhan}
溶度积还能算出一个反直觉的现象：\textbf{同离子效应}。前面算过，氯化银在纯水里的溶解度约为 $s=\sqrt{1.8\times10^{-10}}\approx1.3\times10^{-5}\ \mathrm{mol/L}$。现在把它放进 \ud{0.1}{mol/L} 的食盐水里——水里本来就有大量氯离子。设此时溶解度为 $s'$，则 $[\chem{Ag^+}]=s'$，$[\chem{Cl^-}]\approx0.1$（食盐贡献的氯离子远比氯化银溶出的多），于是
\[K_{\mathrm{sp}} = s'\times0.1 \approx 1.8\times10^{-10} \quad\Rightarrow\quad s'\approx1.8\times10^{-9}\ \mathrm{mol/L}\]
溶解度暴跌到原来的约七千分之一！直觉上“加点盐应该没关系”，账本上却是天差地别。这正是“加入含同种离子的物质会压低溶解度”的定量版，也是化学工业里“盐析”回收产品、实验室里让沉淀更完全的原理。动手算一算：如果食盐水浓度改成 \ud{0.01}{mol/L}，溶解度又是多少？（跳过本节不影响主线。）
\end{tiaozhan}

\section*{练一练}

\begin{sikaoti}
\item 硬水烧开会结水垢，但海水晒干了留下的是盐而不是“垢”。从溶度积的角度解释：为什么氯化钠不会在水壶里析出，碳酸钙却会？（提示：想想两者溶解度的数量级差别。）
\item 某同学检验水样中的氯离子：滴入硝酸银溶液后出现白色沉淀，他立刻宣布“水里有氯离子”。请指出他漏掉了哪一步关键操作，这一步是为了排除哪些“冒牌货”？画出你的检验流程图（取样 → 加试剂 → 观察 → 排除 → 结论）。
\item 画一幅粒子示意图：用圆圈表示铜离子、小尾巴球表示氨分子，画出“加少量氨水出现沉淀”和“加过量氨水沉淀溶解变深蓝”两个阶段的粒子图景，标出中心离子、配体和配位数。图旁注明：这是人为的表示法，真实的离子没有颜色和形状。
\item 钡餐用的是有毒元素钡的化合物，却能安全喝下去。用溶度积的语言写一段话，向一位担心钡餐安全的朋友解释它为什么安全，并说明“元素本身有毒”和“化合物可以服用”为什么不矛盾。
\item 一条河流下游的村庄怀疑上游工厂排污。环保人员到场后要回答两个问题：“水里有没有铅？”和“水里还有什么我们没想到的东西？”分别该选哪类侦探方法（试剂检验、光谱、色谱、质谱）？为什么一个问题不能只靠一种方法？
\item 章末挑战里，氯化银在食盐水中溶解度暴跌。反过来想：为什么把氯化银泡在浓氨水里它反而溶解得更多？用“平衡被搬走”的思路，画一幅“溶解平衡被配位反应拉动”的物质流图。
\end{sikaoti}
